\section{Step 9 Publish the completed loader result}\label{png-load-s09}
Fill \texttt{load\_publish}. Return the candidate in \texttt{\$v0}; initialize \texttt{\$v1} to zero. A failed candidate goes to the supplied restore path with count zero. For success, put the retained byte count in \texttt{\$v1} and continue through \texttt{load\_publish\_ready} to restore.

The restore code reloads the caller's saved registers and stack pointer. Copying the count to its return register before restoring \texttt{\$s2} preserves the result while returning that saved register to its caller-owned value. The complete caller records 115 and then reaches the still-unfilled decoder, so the program's development message remains expected.

Use \texttt{s09-1} and stop before \texttt{load\_publish\_ready}.

\paragraph{Observe} The direct loader returns status zero and count 115 with saved state restored. The caller publishes loaded\_count but leaves decoder results and pixels untouched.

\paragraph{Check} Why copy s2 into v1 before restoring s2?
\begin{itemize}
  \item[A.] The count becomes a path address during close.
  \item[B.] The restore code deletes the file.
  \item[C.] Restoring s2 recovers the caller's previous value.
\end{itemize}
